% swiftui-navigation.tex — NavigationStack + NavigationPath, navigationDestination(for:), value links,
% programmatic push/pop, NavigationSplitView, sheet/fullScreenCover/popover, dismiss, deep links, router.
% Sources (checked 2026-10-06):
%   developer.apple.com/documentation/swiftui/navigationpath (removeLast(_:) only — no removeAll)
%     · /navigationpath/codable (iOS 16; nil unless every element is Codable) ·
%     /view/navigationdestination(item:destination:) (iOS 17) · (ispresented:destination:) (16)
%     · /navigationsplitview/init(preferredcompactcolumn:sidebar:detail:) (17) ·
%     /view/presentationdetents(_:) (16) · /view/presentationcompactadaptation(_:) (16.4) ·
%     /view/interactivedismissdisabled(_:) (15) · /swiftui/dismissaction (15)
%   NavigationLink init(destination:isActive:label:) · init(tag:selection:destination:label:)
%     (deprecated iOS 16: "Use init(value:label:)") · /swiftui/navigationview (deprecated: "Use
%     NavigationStack and NavigationSplitView")
% @labels: area=swiftui kind=api platform=apple topic=ui,architecture discussion=183
% @tags: navigationstack, navigationpath, navigationdestination, navigationlink, navigationsplitview, sheet, fullscreencover, presentationdetents, dismiss, onopenurl, router, deep-links
\documentclass[8pt]{extarticle}
\usepackage{printup-sheet}

\lstdefinelanguage{SwiftSheet}{
  morekeywords={func,let,var,struct,class,final,enum,protocol,extension,return,if,else,
    guard,case,switch,self,some,any,where,init,in,private,static,
    @State,@Binding,@Observable,@Environment,@Bindable,@MainActor},
  alsoletter={@}, sensitive=true, morecomment=[l]{//}, morestring=[b]"}

\begin{document}

\sheettitle{SwiftUI navigation}{swiftui · folio}

\oneliner{Since iOS 16 navigation is \textbf{data}: a \texttt{NavigationStack} renders a \textbf{path} (an
array of \texttt{Hashable} values) that \emph{you} own; \texttt{navigationDestination(for:)} maps each value
\emph{type} to a screen. Push = append, pop = remove, deep link = assign a whole path.}

\begin{multicols}{2}
\raggedright

\section{How it works}
\begin{itemize}
  \item \textbf{\texttt{NavigationStack(path:\,root:)}} (16) — the path is \texttt{[Route]} (one type) or
        \textbf{\texttt{NavigationPath}}: a \emph{type-erased} list of \texttt{Hashable} values, so one stack
        can hold \texttt{Item}, \texttt{User}, \texttt{Int}\ldots, each routed by its own destination.
  \item \textbf{\texttt{NavigationLink(\_:value:)}} pushes a \emph{value}, not a view — lazy, decoupled,
        programmatic. \texttt{NavigationLink \{ Dest() \} label: \{\}} (view-based) still works but cannot be
        driven from the path.
  \item \textbf{\texttt{.navigationDestination(for: T.self) \{ v in … \}}} — \textbf{one} per type per stack,
        inside the stack, on a view that is always there (not in a lazy row).
        iOS 17 adds \texttt{navigationDestination(item:)}; \texttt{(isPresented:)} is 16.
  \item \textbf{Programmatic}: push \texttt{path.append(x)}; pop \texttt{path.removeLast()};
        pop-to-root \texttt{path = []} / \texttt{path.removeAll()} for an array, or
        \texttt{path.removeLast(path.count)} / \texttt{path = NavigationPath()} — \texttt{NavigationPath}
        has \emph{no} \texttt{removeAll}.
  \item \textbf{Restoration}: if every element is \texttt{Codable}, \texttt{path.codable} gives a
        \texttt{NavigationPath.CodableRepresentation} (else \texttt{nil}); encode to \texttt{@SceneStorage}
        / disk, rebuild with \texttt{NavigationPath(rep)}.
  \item \textbf{\texttt{NavigationSplitView}} (16) — 2 columns (\texttt{sidebar}, \texttt{detail}) or 3
        (+\texttt{content}); selection drives the next column; \texttt{columnVisibility} binding;
        on compact width it \textbf{collapses into a stack} (\texttt{preferredCompactColumn}, 17).
  \item \textbf{Modals}: \texttt{.sheet} (card, swipe-dismiss, \texttt{.presentationDetents([.medium,
        .large])} 16); \texttt{.fullScreenCover} (no swipe-dismiss); \texttt{.popover} (bubble on iPad/Mac,
        a sheet on compact unless \texttt{.presentationCompactAdaptation}, 16.4). Each takes
        \texttt{isPresented:} (\texttt{Bool}) or \texttt{item:} (\texttt{Identifiable?} — preferred when
        content depends on data). Block swipe: \texttt{.interactiveDismissDisabled()}.
  \item \textbf{\texttt{@Environment(\textbackslash.dismiss)}} (15) — pops the pushed view \emph{or} closes the
        sheet/cover that presented this view: the \emph{nearest} context.
  \item \textbf{Tabs}: one \texttt{NavigationStack} \emph{per tab}, each with its own path.
  \item \textbf{Router}: an \texttt{@Observable} object owning \texttt{path} (+ sheet state), put in the
        environment, so any view calls \texttt{router.show(.x)} and \texttt{.onOpenURL} maps a URL to routes.
\end{itemize}

\section{Example — router + deep link}
\begin{lstlisting}[language=SwiftSheet]
enum Route: Hashable, Codable { case item(Int), user(String) }
@Observable final class Router {
  var path: [Route] = []
  func open(_ url: URL) { path = Route.parse(url) } // all at once
}
struct Root: View {
  @State private var router = Router()
  var body: some View {
    NavigationStack(path: $router.path) {
      List(1..<50, id: \.self) { n in
        NavigationLink("Item \(n)", value: Route.item(n)) }
      .navigationDestination(for: Route.self) { Screen(route: $0) }
    }
    .environment(router).onOpenURL { router.open($0) }
  } }
\end{lstlisting}

\columnbreak

\section{The stack \emph{is} the path}
\begin{tikzpicture}[sheet]
  \tikzset{pv/.style={cell, minimum width=15mm, minimum height=6mm, fill=sheetBlue!8, draw=sheetBlue},
           sc/.style={box, minimum width=15mm, minimum height=9mm, font=\footnotesize, draw=sheetGreen, fill=sheetGreen!8}}
  \node[font=\bfseries\small, anchor=west] at (-0.75,0.6) {\texttt{path} = the data you own};
  \node[pv, fill=black!5, draw=sheetGrey] (r)  at (0,0)   {(root)};
  \node[pv] (p0) at (1.7,0)  {\texttt{.item(3)}};
  \node[pv] (p1) at (3.4,0)  {\texttt{.user("a")}};
  \node[pv] (p2) at (5.1,0)  {\texttt{.item(7)}};
  \node[cell, dashed, draw=sheetOrange, minimum width=9mm, minimum height=6mm, text=sheetOrange] (slot) at (6.6,0) {+};
  \draw[hot] (p2.north) to[bend left=40] node[above, note, text=sheetOrange]{\texttt{append} = push} (slot.north);
  \node[sc, fill=black!5, draw=sheetGrey] (s)  at (0,-1.7)   {List};
  \node[sc] (s0) at (1.7,-1.7) {Detail 3};
  \node[sc] (s1) at (3.4,-1.7) {Profile a};
  \node[sc, very thick] (s2) at (5.1,-1.7) {Detail 7\\[-2pt]\scriptsize\itshape on top};
  \foreach \a/\b in {r/s,p0/s0,p1/s1,p2/s2} \draw[flow] (\a) -- (\b);
  \node[note, align=left, anchor=west] at (5.95,-1.7) {\scriptsize screens\\[-2pt]\scriptsize SwiftUI\\[-2pt]\scriptsize renders};
  \node[note, align=left, anchor=west] at (5.95,-0.82) {\scriptsize one\\[-2pt]\scriptsize \texttt{destination}\\[-2pt]\scriptsize per type};
  % operations
  \draw[->, thick, sheetRed] (4.85,-2.15) -- (4.85,-2.45) -- node[below, note, text=sheetRed]{\texttt{removeLast()} = pop} (3.4,-2.45) -- (3.4,-2.15);
  \draw[->, thick, sheetRed, dashed] (5.35,-2.15) -- (5.35,-3.05) -- node[below, note, text=sheetRed, pos=0.55]{\texttt{path = []} = pop to root} (0,-3.05) -- (0,-2.15);
  % deep link lane
  \node[box, draw=sheetBrown, fill=sheetBrown!8, font=\footnotesize, anchor=west] (u) at (-0.75,-4.1) {URL \texttt{app://item/3/user/a}};
  \node[box, draw=sheetBlue, font=\footnotesize, anchor=east] (pp) at (7.1,-4.1) {\texttt{path = [.item(3), .user("a")]}};
  \draw[hot] (u) -- node[above, note, text=sheetOrange]{parse} (pp);
  \node[note, anchor=west, text width=77mm] at (-0.75,-4.8)
       {Deep link = \emph{assign} the whole array; SwiftUI pushes every screen, back button included.
        Codable routes $\Rightarrow$ the same array restores state on relaunch.};
\end{tikzpicture}

\section{Traps}
\begin{itemize}
  \trap{\texttt{NavigationLink(isActive:)} / \texttt{(tag:selection:)} are \textbf{deprecated} (iOS 16), and so
        is \texttt{NavigationView}: one \texttt{Bool} per link, buggy pops, no restoration.}
  \trap{\texttt{navigationDestination} inside a \texttt{List}/\texttt{LazyVStack} row, or two for one type
        — ``no matching destination'' or the wrong screen. Declare once, on a stable ancestor.}
  \trap{\texttt{.sheet(isPresented:)} + separate \texttt{@State var selected} — can present stale/\texttt{nil}
        data. Use \texttt{.sheet(item: \$selected)}.}
  \trap{\texttt{dismiss()} inside a \texttt{NavigationStack} \emph{within} a sheet pops the inner stack,
        it does not close the sheet — pass a closure/binding or read \texttt{dismiss} at sheet level.}
  \trap{\texttt{.navigationTitle} goes on the \emph{content}, not on \texttt{NavigationStack}; a split view
        tested only on iPad shows a blank detail when collapsed on iPhone.}
  \trap{Nesting a \texttt{NavigationStack} in a pushed screen = two bars / broken back stack.}
\end{itemize}

\section{Remember}
\textbf{``Push values, not views.''} Path = \texttt{[Hashable]}; destination = \emph{per type};
\texttt{item:} beats \texttt{isPresented:}; \texttt{dismiss} = nearest context.


\end{multicols}

\noindent{\footnotesize\color{sheetGrey}\textit{Related:} ios-swift/navigation-presentation (UIKit) ·
\texttt{TabView} / \texttt{Tab} (18) · \texttt{@SceneStorage} · universal links · \texttt{.toolbar} ·
\texttt{presentationDetents}}

\end{document}
